\label{sec:problem-setup}

\paragraph{Road-graph route generation.}
We study conditional \emph{route generation} at the trip-segment level.
We represent the environment as a directed road graph $G=(V,E)$ extracted from OpenStreetMap and discretized into a shared grid coordinate system.
Each node $v\in V$ has a grid coordinate $(y_v,x_v)\in[0,1023]^2$; each directed edge $e=(u\!\rightarrow\!v)\in E$ has a length and a road-tier type (major/minor/service).
A ground-truth trip segment is represented as a node sequence
$y=(v_0,\dots,v_L)$ with $v_0=o$ and $v_L=d$, obtained by map-matching a continuous trajectory to $G$.

\paragraph{Conditioning.}
Each sample is conditioned on
$c=(o,d,t_0,x)$,
including origin $o$, destination $d$, departure time $t_0$ (Unix epoch seconds), and optional context $x$.
In this paper we focus on a minimal context set for corridor choice: cyclical time features derived from $t_0$ and road-tier information from $G$.

\paragraph{Goal.}
Given $c$, the model samples corridor paths $\hat{y}\sim p_\theta(y\mid c)$ such that:
(i) \textbf{intent consistency} (reaching $d$),
(ii) \textbf{graph feasibility} (the output is a valid path on $G$), and
(iii) \textbf{corridor multi-modality} (supporting distinct topological alternatives).

\paragraph{Why continuous end-to-end generation is brittle.}
Even for fixed $(o,d,t_0)$, multiple corridor-level alternatives can be far apart in coordinate space while remaining intent-equivalent.
This multi-modality is structurally discrete: the main uncertainty lies in \emph{which corridor} to take, while within-corridor variation is comparatively continuous.
Representing corridor choice only implicitly in a continuous coordinate output space leads to destructive averaging or long-horizon drift, motivating explicit \emph{corridor commitment} on the road graph.

\paragraph{Evaluation: best-of-$K$ corridor match and success.}
Given a ground-truth path $y$ and $K$ sampled paths $\{\hat{y}^{(k)}\}_{k=1}^K$, we evaluate corridor-level match by edge-set Jaccard similarity.
Let $E(y)$ be the set of directed edges traversed by $y$.
We define
\begin{equation}
  J(\hat{y},y)=\frac{|E(\hat{y})\cap E(y)|}{|E(\hat{y})\cup E(y)|},
  \quad
  J_{\mathrm{best}}=\max_{k\in\{1,\dots,K\}} J(\hat{y}^{(k)},y).
\end{equation}
We report the mean/median/90th percentile of $J_{\mathrm{best}}$ across routes.
We also report a \textbf{success rate}, defined as the fraction of samples for which the executor returns a complete path from $o$ to $d$ (e.g., A* succeeds for all waypoint-to-waypoint segments) within a maximum step budget.

\paragraph{Multi-modal OD groups for gates.}
To test whether context carries measurable signal for corridor-level variability, we group routes into coarse OD buckets (by quantizing $(o,d)$ into OD-bin keys) and cluster the ground-truth graph paths within each bucket using an edge-Jaccard distance threshold.
This yields corridor-cluster labels for a subset of sufficiently multi-modal OD buckets, enabling a lightweight gate that quantifies the predictive signal of $(t_0,\text{tier})$ via AUC (Sec.~\ref{sec:experiments}).
